\subsubsection{Actitudes y experiencias: contraste con la encuesta de 2017}
La ENVCM 2017 amplía el contexto con experiencias de mujeres de 15 años o
más. Sus porcentajes nacionales de violencia en ámbitos público y privado
son 67,4\,\% a lo largo de la vida y 33,8\,\% en los doce meses previos
(cuadro \ref{tab:envcm}). Incluyen ámbitos distintos del hogar; 67,4\,\% no
es una prevalencia exclusiva de violencia doméstica \citep{envcm2017}.

\begin{table}[htbp]\centering\small
\caption{Violencia experimentada por mujeres: contexto nacional de 2017}\label{tab:envcm}
\begin{tabularx}{\textwidth}{@{}Yrrp{2.9cm}@{}}\toprule
Período de referencia & \% & IC95 publicado & Universo\\\midrule
A lo largo de la vida & 67,4 & 65,2--69,6 & Mujeres de 15 años o más\\
Últimos doce meses & 33,8 & 31,7--36,0 & Mujeres de 15 años o más\\\bottomrule
\end{tabularx}
\fuente{DIGESTYC (2018), ENVCM 2017, tablas 46--47, pp. 83--84, fila ``Ámbito público y privado''. Puntos e IC transcritos del informe, no estimados por este artículo.}
\end{table}

La diferencia entre ambos períodos no indica una caída temporal: una ventana
anual recoge menos exposición que una historia de vida. Tampoco se resta
7,7\,\% de justificación MICS a estos porcentajes para medir contradicción
o subregistro: cambian año, edades, informantes y constructo. Su aportación
es mostrar por qué experiencias, normas y delitos registrados requieren
instrumentos complementarios, con denominadores explícitos.
